\documentclass[aspectratio=169,11pt]{beamer}

\usepackage[T1]{fontenc}
\usepackage[utf8]{inputenc}
\usepackage{lmodern}
\usepackage{amsmath,amssymb}
\usepackage{booktabs}
\usepackage{xcolor}
\usepackage{hyperref}
\usepackage{microtype}

\usetheme{Madrid}
\usecolortheme{default}

\definecolor{LaurierPurple}{HTML}{4B116F}
\definecolor{DeepBlue}{HTML}{0047AB}

\setbeamertemplate{navigation symbols}{}
\setbeamercolor{title}{fg=white,bg=LaurierPurple}
\setbeamercolor{frametitle}{fg=white,bg=LaurierPurple}
\setbeamercolor{block title}{fg=white,bg=DeepBlue}
\setbeamercolor{structure}{fg=DeepBlue}

\hypersetup{colorlinks=true,urlcolor=DeepBlue,linkcolor=DeepBlue}

\title[Replication]{Module 8: Empirical Replication}
\subtitle{EC 410K / EC 610I --- Machine Learning in Economics}
\author{M. Jahangir Alam, PhD}
\institute{Wilfrid Laurier University}
\date{Spring 2026}

\begin{document}

\begin{frame}[plain]
  \titlepage
\end{frame}

\begin{frame}{Module overview}
  \begin{itemize}
    \item Replication is the bridge between reading papers and doing research.
    \item Goals: verify computational claims, understand design details, and create a foundation for extensions.
    \item Outputs: clean code, documented data, and a short memo explaining matches and discrepancies.
  \end{itemize}
\end{frame}

\begin{frame}{Learning objectives}
  \begin{enumerate}
    \item Build a replication checklist from a paper’s empirical section.
    \item Diagnose common failures: sample mismatch, wrong clustering, missing weights, version drift.
    \item Practice writing reproducible workflows suitable for GitHub + Colab.
    \item Connect replication to the course’s ML extension project.
  \end{enumerate}
\end{frame}

\begin{frame}{What counts as a successful replication?}
  \begin{itemize}
    \item Reproduce \textbf{one} core exhibit well (table, figure, key coefficient).
    \item Document every step: data pull, cleaning, estimation command, package versions.
    \item If you cannot match: explain the most likely reasons and what you tried.
  \end{itemize}
\end{frame}

\begin{frame}{Replication and identification}
  \begin{itemize}
    \item Replication clarifies \textbf{what variation} the estimate uses.
    \item It does not by itself validate assumptions---but it makes scrutiny possible.
  \end{itemize}
\end{frame}

\begin{frame}{Course replication papers (reminder)}
  \begin{itemize}
    \item Acemoglu--Johnson--Robinson (2001); Nunn (2008); Goldin--Lurie--McCubbin (2016);
    \item Hansen (2015); Anderson (2017); Autor--Dorn--Hanson (2013).
  \end{itemize}
\end{frame}

\begin{frame}{Workflow (recommended)}
  \begin{enumerate}
    \item Read the empirical strategy before coding.
    \item Obtain data; verify units and merges.
    \item Reconstruct the baseline sample.
    \item Match a simple exhibit first; then complex robustness.
    \item Version control + environment snapshot (\texttt{requirements.txt} / conda).
  \end{enumerate}
\end{frame}

\begin{frame}{Python / Colab demonstration}
  \begin{itemize}
    \item Toy DGP where the target parameter is known.
    \item Compare OLS estimate to a fake ``published'' value to practice writing a replication memo.
  \end{itemize}
\end{frame}

\begin{frame}[fragile]{Colab setup}
  \begin{verbatim}
!pip -q install numpy pandas statsmodels
  \end{verbatim}
\end{frame}

\begin{frame}{Student / group assignment}
  \begin{itemize}
    \item Create a two-page replication memo draft for your assigned paper: question, data sources, baseline spec, target exhibit.
    \item Run the toy Colab exercise and use it as a template for discussing mismatch sources.
    \item Present progress: what is reproduced, what is blocked, next steps toward ML extension.
  \end{itemize}
\end{frame}

\begin{frame}{Summary}
  \begin{itemize}
    \item Replication is a disciplined craft: transparency, testing, documentation.
    \item Treat mismatch as information, not failure---but require a clear debugging plan.
    \item Strong replication makes the ML extension scientifically meaningful.
  \end{itemize}
\end{frame}

\begin{frame}[plain]
  \begin{center}
    \Large Questions?\\[0.8em]
    \normalsize Next: applied ML + causal project integration (Module 9).
  \end{center}
\end{frame}

\end{document}
